\section{Introduction}

Software has become an integral and integrated part of everyday life, accompanied
by a near constant stream of updates to various devices all around
us. Unfortunately, software updates are a frequent source of issues
\autocites{zhangUnderstandingDetectingSoftware2021}[][]{Gray1986WhyDC}, often
making it into production environments before the necessary edge cases are
triggered. We believe that techniques from formal verification can be applied
here to ensure update compatibility, which is ultimately what my project aims to
do. Unlike some previous work
\autocites{ajmaniModularSoftwareUpgrades2006}[][]{reitblattAbstractionsNetworkUpdate2012},
I aim to prove compatibility of software updates without imposing restrictions
about how the update itself is performed.

The most nebulous part of this goal is how to define ``compatible'', which I
break down into two major categories; data compatibility and operational
compatibility. Data compatibility aims to show that that the new version can
correctly interpret persistent state left behind by the pervious version while
operational compatibility reasons about behavior of invoking the same function
or feature in both versions of the software. This project is currently working
on data compatibility.

There are numerous data description languages, often paired with a particular
encoding format. Notable examples of this include ASN1
\cite{ASN1EncodingRules2021}, protobuf \cite{LanguageGuideProto}, CBOR
\cite{birkholzConciseDataDefinition2019,bormannConciseBinaryObject2020} and even
JSON with JSONschema
\cite{wrightJSONSchemaMedia2022,brayJavaScriptObjectNotation2017} or using
structs in some host programming language. All of these formats have seen prior
formalization in proof-assistants or proof oriented programming languages
\cite{habibFindingDataCompatibility2021,ramananandroSecureParsingSerializing2025,yeVerifiedProtocolBuffer2019,niASN1ProvablyCorrect2023},
although none of these works consider the common real-world case of an update to
the schema used to pass messages between nodes of a distributed system or store
data while the software isn't running. Some tools already exist which can report
on if two versions of a schema file are compatible (i.e. Buf), although these
tools operate using ad-hoc notions of compatibility without formalizing a
specification or verifying the implementation.

Using Lean, I plan to verify a compatibility checker for one or more of
these libraries before expanding the project to consider a wider range of data
formats. It should also be possible to include some analysis of the source code
operating on the state to prove more invasive updates, or even the proofs of
verified pieces of software, although that is a goal for a longer timescale.

\subsection{Versioned Formats}

Now, technically speaking, this problem is solved from a practical
standpoint. What you'd do is encode the version number at the very beginning of
the encoded blob and then parse just that data at first in order to decide which
versions parser to use. However, many commonly used data descriptor languages
like Protobuf, JSON, ASN1, etc, don't have built in mechanisms for this. Some
encode which version of the wire protocol they use, which is a helpful feature,
but not which version of \emph{your message} this is. For large messages, having
to parse under multiple descriptors might be an unacceptable performance hit.

\subsection{Impact on Verified Systems}

With the rise of AI and agentic proofs, expect to see more and more verified
systems implemented and deployed. As more verified systems enter the real world,
eventually they'll need to be updated. However, the verification process will
likely only verify that if a message is encoded under the correct descriptor,
the system will behave correctly. Encountering a message encoded under an old
descriptor side-steps the verification guarantees entirely.

% Local Variables:
% citar-bibliography: ("../pollux.bib")
% TeX-command-default: "Make"
% TeX-master: "../pollux.tex"
% End:
